\documentclass[12pt]{article}
\usepackage{amsmath}
\begin{document}

\textit{Proposed Solution to \#5816 SSMJ }\\ \\
\textit{Solution proposed by Prakash Pant, The University of Vermont, Bardiya, Nepal.}\\ \\
\textit{Problem Proposed by Daniel Sitaru, National Economic College "Theodor Costescu"
Drobeta Turnu - Severin, Romania.
} \\ \\
\textbf{Statement of the Problem:}
For \(0 < a \le b \), prove:\\

\[
\int_{a}^{b} \frac{x e^{x^2} \sqrt{e^{x^2}}}{e^{3x^2} + 1} \, dx
\le
\frac{1}{2} \tan^{-1}\!\left(
\frac{e^{b^2} - e^{a^2}}{1 + e^{a^2 + b^2}}
\right).
\]
\textbf{Solution of the Problem:}
Observe that : 
\[ \frac{1}{2} \tan^{-1}\!\left(
\frac{e^{b^2} - e^{a^2}}{1 + e^{a^2 + b^2}}
\right) = \frac{1}{2} \left( \tan^{-1}(e^{b^2}) - \tan^{-1}(e^{a^2}) \right) = \frac{1}{2} ( \tan^{-1}(e^{x^2})) |_a^b \]
\[  =\frac{1}{2} \int_a^b \frac{d}{dx}(\tan^{-1}(e^{x^2})) dx  =  \int_a^b \frac{xe^{x^2}}{1+e^{2x^2}} dx  \]
Now, the problem reduces to proving 
\[
\int_{a}^{b} \frac{x e^{x^2} \sqrt{e^{x^2}}}{e^{3x^2} + 1} \, dx
\le
 \int_a^b \frac{xe^{x^2}}{1+e^{2x^2}} dx 
\]
for which it suffices to prove
\[ \frac{\sqrt{e^{x^2}}}{e^{3x^2}+1} \le \frac{1}{1+e^{2x^2}} \]
Since the denominators are clearly positive, we can cross multiply
\begin{equation} e^{\frac{x^2}{2}} + e^{\frac{5x^2}{2}} \le e^{3x^2}+1
\end{equation}
To prove this statement, we use Karamata Inequality.\\ Since sequence \( (3,0) \succ (\frac{5}{2},\frac{1}{2})\) and \( f(a) = e^{ax^2}\) is convex whenever a is positive, thus by Karamata inequality, 
\[ f(3) + f(0) \ge f(\frac{5}{2}) + f(\frac{1}{2})\]
Using \( f(a) = e^{ax^2}\) proves equation (1) and hence the original inequality. 
\end{document}